\chapter{Technology Stack Comparison and Architectural Justification}
\label{chap:tech_justification}

\section{Frontend Framework Selection}
The choice of frontend framework affects developer productivity, rendering performance, and long-term maintainability. The Lectura client needs to handle complex document rendering, real-time AI generation updates, and interactive components like editable slides. Table~\ref{tab:frontend_comparison} compares the main JavaScript/TypeScript frameworks considered.

\begin{table}[htbp]
\centering
\resizebox{\textwidth}{!}{%
\begin{tabular}{|l|p{4cm}|p{4cm}|p{4cm}|}
\hline
\textbf{Criterion} & \textbf{React (Selected)} & \textbf{Vue.js} & \textbf{Angular} \\ \hline
State Management & Flexible (Zustand, Redux, React Query). & Built-in (Pinia). & Standardized (RxJS). \\ \hline
Ecosystem & Very large third-party library ecosystem. & Growing but smaller. & Enterprise-focused. \\ \hline
Rendering & Virtual DOM reconciliation. & Virtual DOM reconciliation. & Incremental DOM updates. \\ \hline
TypeScript Support & Full TypeScript integration. & Excellent native support. & TypeScript required. \\ \hline
\end{tabular}%
}
\vspace{6pt}
\caption{Comparison of frontend SPA frameworks.}
\label{tab:frontend_comparison}
\end{table}

\textbf{Justification:} React 18 was selected primarily for its modular component architecture and mature ecosystem. The project required rendering dynamic content types — including mathematical LaTeX notation in quizzes and generated Markdown tables — and React's declarative model makes it straightforward to integrate specialized rendering libraries without complicating the main application logic. React Query was used to handle server-state caching, which prevents unnecessary network requests when a user switches between different generated summaries and flashcard decks.

\section{Backend Language Selection}
The backend needs to handle concurrent workloads efficiently: processing large JSON payloads, maintaining active WebSocket connections, and coordinating long-running AI generation tasks. Table~\ref{tab:backend_comparison} compares the three main candidates.

\begin{table}[htbp]
\centering
\resizebox{\textwidth}{!}{%
\begin{tabular}{|l|p{4cm}|p{4cm}|p{4cm}|}
\hline
\textbf{Criterion} & \textbf{Go (Selected)} & \textbf{Node.js (TypeScript)} & \textbf{Python (FastAPI)} \\ \hline
Concurrency Model & Native goroutines with precise memory control. & Single-threaded event loop. & Threading / AsyncIO (limited by GIL). \\ \hline
Execution Speed & Very high (compiled to native binary). & High (JIT compilation via V8). & Moderate (interpreted). \\ \hline
Memory Usage & Low. & Moderate to high. & High. \\ \hline
Type Safety & Strict, enforced at compile time. & Optional (TypeScript). & Runtime type hints only. \\ \hline
\end{tabular}%
}
\vspace{6pt}
\caption{Backend language comparison for concurrent workload handling.}
\label{tab:backend_comparison}
\end{table}

\textbf{Justification:} Go was chosen because the Lectura architecture relies heavily on a Worker Pool pattern to handle LLM generation tasks that can take 10--60 seconds each. Go's goroutines and channel-based communication provide an efficient way to manage many concurrent tasks with low memory overhead. Node.js's single-threaded event loop can struggle under heavy computational load, and Python's Global Interpreter Lock (GIL) limits true parallel execution on multi-core machines. Go compiles to a single binary with fast startup times and low memory usage, which also helps reduce cloud hosting costs.

\section{Database Selection}
Choosing a database required balancing two different needs: structured relational data for user accounts and authentication, and flexible schema storage for the unpredictable JSON output from AI generation. Table~\ref{tab:database_comparison} compares the options.

\begin{table}[htbp]
\centering
\resizebox{\textwidth}{!}{%
\begin{tabular}{|l|p{4.5cm}|p{3.5cm}|p{3.5cm}|}
\hline
\textbf{Criterion} & \textbf{PostgreSQL (Selected)} & \textbf{MongoDB} & \textbf{MySQL} \\ \hline
Data Integrity & Full ACID compliance. & Eventually consistent. & ACID compliant. \\ \hline
Schema Flexibility & Strong (JSONB with indexing). & Native document store. & Limited. \\ \hline
Query Capabilities & Advanced (array ops, CTEs, window functions). & Good for document queries. & Standard SQL. \\ \hline
Complex Joins & Well optimized. & Poor performance. & Well optimized. \\ \hline
\end{tabular}%
}
\vspace{6pt}
\caption{Database comparison based on schema flexibility and query performance.}
\label{tab:database_comparison}
\end{table}

\textbf{Justification:} PostgreSQL was chosen because it handles both use cases well. While MongoDB is better suited for purely unstructured data, it lacks the relational guarantees needed for user authentication and progress tracking. PostgreSQL avoids the need for two separate databases by combining standard relational tables with its JSONB column type. This allows the backend to store dynamically sized arrays of generated quiz questions or summary data in a single column while still maintaining proper foreign key relationships and data normalization for user-related tables.

\section{LLM Provider Selection}
The core intelligence of Lectura depends on the external Large Language Model used for content generation. The key factors are context window size, response speed, and cost per token. Table~\ref{tab:llm_comparison} compares the main options.

\begin{table}[htbp]
\centering
\resizebox{\textwidth}{!}{%
\begin{tabular}{|l|p{4cm}|p{4cm}|p{4cm}|}
\hline
\textbf{Criterion} & \textbf{Google Gemini 3 Flash} & \textbf{OpenAI GPT-4o} & \textbf{Anthropic Claude 3.5} \\ \hline
Context Window & 1+ million tokens. & 128,000 tokens. & 200,000 tokens. \\ \hline
Response Speed & Very fast. & Moderate. & Moderate. \\ \hline
Cost & Low. & High. & High. \\ \hline
Format Compliance & Good. & Excellent. & Very good. \\ \hline
\end{tabular}%
}
\vspace{6pt}
\caption{Comparison of LLM providers by context window, speed, and cost.}
\label{tab:llm_comparison}
\end{table}

\textbf{Justification:} Google Gemini 3 Flash was selected primarily for its large context window. In an educational application where the input might be a full lecture transcript or even an entire textbook chapter, smaller context windows (like GPT-4o's 128K tokens) would require implementing a Retrieval-Augmented Generation (RAG) pipeline with vector embeddings — adding significant architectural complexity. Gemini's million-token context window allows the system to pass the entire source document in a single API call, which simplifies the architecture and ensures that no content is lost during chunking. Its lower cost per token and fast response times also make it practical for a student-facing application where usage volume can be high.

\section{Document Export Approach}
Exporting generated web content (including CSS layouts, mathematical formulas, and styled presentation slides) into downloadable PDF and PPTX files is a non-trivial technical challenge. Table~\ref{tab:export_comparison} compares the available approaches.

\begin{table}[htbp]
\centering
\resizebox{\textwidth}{!}{%
\begin{tabular}{|l|p{4cm}|p{4cm}|p{4cm}|}
\hline
\textbf{Criterion} & \textbf{md-to-pdf (Puppeteer)} & \textbf{jsPDF} & \textbf{pdfmake} \\ \hline
Execution & Server-side headless browser. & Client-side Canvas API. & Client-side JSON-based. \\ \hline
Layout Support & Full CSS/HTML rendering. & Manual coordinate positioning. & Abstract layout definitions. \\ \hline
Responsive Scaling & Consistent output regardless of viewport. & Viewport-dependent. & Not applicable. \\ \hline
Math Rendering & Full support via browser engine. & Very limited. & Very limited. \\ \hline
\end{tabular}%
}
\vspace{6pt}
\caption{Comparison of PDF export methods.}
\label{tab:export_comparison}
\end{table}

\textbf{Justification:} Client-side libraries like jsPDF were not suitable because they struggle with complex layouts and mathematical notation. For presentations specifically, client-side exports can produce distorted results when the user's browser viewport differs from the intended slide dimensions ($1161 \times 653$ pixels). Server-side rendering using a headless Chromium instance (via Puppeteer) avoids these issues by rendering the content in a controlled environment. The system constructs the DOM in a headless browser, renders it at the correct dimensions, and captures a pixel-accurate PDF. This ensures that exported study materials look identical to what users see in the web interface.

\subsection{Memory Management for Server-Side Rendering}
While server-side rendering with Puppeteer produces high-quality exports, each headless Chromium instance consumes a significant amount of memory. Launching a new browser process for every export request would quickly exhaust the server's RAM.

To manage this, the backend maintains a pool of pre-started Chromium instances. Each instance is constrained within defined memory limits using container-level resource controls. The total memory usage $M_{\text{total}}$ for $n$ concurrent exports is bounded by:

\begin{equation}
    M_{\text{total}}(n) = M_{\text{base}} + n \times (M_{\text{instance}} + M_{\text{content}})
\end{equation}

Where $M_{\text{base}}$ is the fixed overhead of the main service and $M_{\text{instance}}$ is the memory required for the Chromium binary. To prevent memory leaks — a known issue where Chromium does not fully release memory after rendering complex pages — the backend tracks how many documents each instance has processed. After reaching a set threshold of $k$ documents, the instance is terminated and replaced with a fresh one. This recycling approach keeps memory usage predictable even during high-traffic periods like exam weeks.